\ifdefined\iscombined
	\lecturetitle{Equivalence Classes}
\else
\documentclass{article}
\usepackage{listings}
\usepackage{amsthm}
\usepackage{comment}
\usepackage{amsmath}
\usepackage{amsfonts}
\usepackage{sectsty}
\usepackage{hyperref}
\usepackage{fancyhdr}
\usepackage[top=1in,bottom=1in,left=1in,right=1in]{geometry}
\usepackage{float}
\usepackage{graphicx}
\usepackage{tabularx}
\usepackage{mathtools}

\date{
	\vspace{-.75em}
	{\large \today}
}

\title{
	\vspace*{-1.5em}
	{\Huge Lecture 9} \\
	\vspace{-1em}
	\rule{\textwidth/4}{.01em} \\
	\vspace{-.25em}
	{\Large Equivalence Classes}
	\vspace{-.75em}
}

\author{
	{\large Matthew Farah}
}

\hypersetup{
	colorlinks=true,
	linkcolor=blue,
	filecolor=magenta,
	urlcolor=blue,
	citecolor=blue,
	anchorcolor=blue
}

\fancyhead{}
\fancyhead[L]{\today}
\fancyhead[R]{Lecture 9 - Equivalence Classes}

\sectionfont{\fontsize{12}{12}\bfseries}
\subsectionfont{\fontsize{10}{12}\normalfont}
\subsubsectionfont{\fontsize{10}{12}\normalfont}

\begin{document}

\maketitle
\pagestyle{fancy}

\fi

\begin{itemize}
	\item A system's domain is the set of all its possible inputs.
	\item An equivalence class is an equivalence relation for some subset of the program's domain ($E \subseteq D$).
	\begin{itemize}
		\item For $x \in D$, $E_x$ is an equivalence class for $x$ if $\forall y \in D; x \equiv y \implies y \in E_x$.
	\end{itemize}
	\item If two inputs belong to the same equivalence class, the program behaves the same for both inputs.
	\begin{itemize}
		\item Useful to reduce the number of tests needed.
	\end{itemize}
	\item A system's domain can be `partitioned' into multiple equivalence classes.
	\item $P = \set{E_{x_{1}}, E_{x_{2}}, \ldots, E_{x_{n}}}$ is a partition of $D$ if its equivalence classes:
	\begin{enumerate}
		\item Are disjoint.
		\begin{itemize}
			\item $\forall i \neq j; E_{x_{i}} \cap E_{x_{j}} = \emptyset$.
		\end{itemize}
		\item Equal the domain.
		\begin{itemize}
			\item $\bigcup_{i = 1}^{n} E_{x_{i}} = D$.
		\end{itemize}
	\end{enumerate}
	\item If the domain can be partitioned, only $n$-many tests (the number of equivalence classes) are needed to test the system.
	\item In practice, very few program domains can be partitioned into meaningful equivalence classes.
	\begin{itemize}
		\item To partition most domains, an equivalence tolerance must be introduced.
		\begin{itemize}
			\item Ex. $x$ is equivalent to $y$ if $|x - y| \le 1$.
			\begin{itemize}
				\item Here, our tolerance is 1.
				\item Tolerance is also called delta.
			\end{itemize}
		\end{itemize}
	\end{itemize}
	\item Partitions with tolerance no longer guarantee that two inputs in the same equivalence class make the system act the same on those inputs.
	\begin{itemize}
		\item Since non-trivial tolerances violate transitivity.
		\item The smaller the tolerance, the less significant this becomes.
	\end{itemize}
	\item A partition $P_1$ finer than partition $P_2$ if $P_2 \subseteq P_1$
	\begin{itemize}
		\item Occurs when $P_1$ has a smaller tolerance than $P_2$.
		\item $P_2$ would be called a courser partition.
	\end{itemize}
	\item A test suite for some partition is called adequate if every equivalence class has a test for it.
	\begin{itemize}
		\item If a test suit is adequate for a partition, it is adequate for courser partitions.
	\end{itemize}
\end{itemize}

\ifdefined\iscombined
	\newpage
\else
	\end{document}
\fi
